\newlength{\abbrleft}
\setlength{\abbrleft}{0.34\textwidth}

\newcommand{\abbr}[2]{%
  \par\noindent
  \parbox[t]{\abbrleft}{\raggedright\strut #1\nobreakspace\textemdash}%
  \hspace{0.02\textwidth}%
  \parbox[t]{\dimexpr\textwidth-\abbrleft-0.02\textwidth\relax}{\strut #2\strut}%
  \par\vspace{0.35\baselineskip}
}

\noindent\textbf{Сокращения}

\medskip

\abbr{ИМ}{имитационное моделирование.}
\abbr{ЛСТ}{преобразование Лапласа--Стилтьеса.}
\abbr{МП}{марковский процесс.}
\abbr{РСМО}{ресурсная система массового обслуживания.}
\abbr{СВ}{случайная величина.}
\abbr{СМО}{система массового обслуживания.}
\abbr{ФР}{функция распределения.}
\abbr{ЦМ}{цепь Маркова.}
\abbr{ЦПТ}{центральная предельная теорема.}

\abbr{BMAP}{\textit{Batch Markovian Arrival Process}, групповой марковский поток событий.}
\abbr{CUDA}{\textit{Compute Unified Device Architecture}, программно-аппаратная архитектура параллельных вычислений.}
\abbr{eMBB}{\textit{Enhanced Mobile Broadband}, класс трафика в сетях 5G для сверхширокополосной связи.}
\abbr{GPU}{\textit{Graphics Processing Unit}, графический процессор.}
\abbr{KL}{дивергенция Кульбака--Лейблера.}
\abbr{MAP}{\textit{Markovian Arrival Process}, марковский поток событий.}
\abbr{MArP}{альтернативное обозначение марковского потока событий (для устранения неоднозначности термина MAP).}
\abbr{ME}{\textit{Matrix-Exponential distribution}, матрично-экспоненциальное распределение.}
\abbr{MMAP}{\textit{Marked Markovian Arrival Process}, маркированный марковский поток событий.}
\abbr{MMPP}{\textit{Markov Modulated Poisson Process}, марковски-модулированный пуассоновский поток.}
\abbr{mMTC}{\textit{Massive Machine Type Communications}, класс массового межмашинного взаимодействия (IoT).}
\abbr{PASTA}{\textit{Poisson Arrivals See Time Averages}.}
\abbr{PH}{\textit{Phase-Type distribution}, распределение фазового типа.}
\abbr{QBD}{\textit{Quasi-Birth-and-Death process}, квазипроцесс рождения и гибели.}
\abbr{TV}{расстояние полной вариации (\textit{total variation distance}).}
\abbr{URLLC}{\textit{Ultra-Reliable Low-Latency Communication}, класс трафика в сетях 5G, критичный к задержкам.}

\bigskip
\noindent\textbf{Основные обозначения}

\medskip

\abbr{\(K\)}{ёмкость системы по числу заявок в редуцированной одноресурсной или одноканальной постановке.}
\abbr{\(N\)}{число обслуживающих приборов; максимальное число заявок, одновременно находящихся в обслуживании.}
\abbr{\(n, k\)}{индекс состояния; число заявок в системе.}
\abbr{\(M\)}{число типов ресурсов.}
\abbr{\(L\)}{число типов заявок.}
\abbr{\(R_m\)}{общий объём ресурса типа \(m\).}
\abbr{\(R\)}{совокупный объём доступного ресурса; в многоресурсной постановке --- совокупность бюджетов ресурсов.}
\abbr{\(r_j = (r_{j1}, \dots, r_{jM})\)}{вектор требований \(j\)-й заявки к ресурсам.}

\abbr{\(A\)}{случайная величина межприходного интервала заявок.}
\abbr{\(W\)}{случайная величина объёма работы (workload) заявки.}
\abbr{\(D\)}{случайная величина ресурсного требования заявки.}
\abbr{\(F(x)\)}{функция распределения требований заявки к ресурсу.}
\abbr{\(F_l(x)\)}{функция распределения требований к ресурсам для заявок типа \(l\).}
\abbr{\(F^{(k)}(x)\)}{\(k\)-кратная свёртка функции распределения \(F(x)\).}
\abbr{\(A_l(x)\)}{функция распределения межприходных интервалов заявок типа \(l\).}
\abbr{\(B(x)\)}{функция распределения объёма работы или времени обслуживания заявки.}
\abbr{\(B_l(x)\)}{функция распределения времени обслуживания заявок типа \(l\).}
\abbr{\(\widetilde{B}(x)\)}{распределение остаточного объёма работы.}
\abbr{\(\mathbb{E}[W], \mathbb{E}[D]\)}{математическое ожидание объема работы и ресурсного требования соответственно.}

\abbr{\(c_v\)}{коэффициент вариации случайной величины.}
\abbr{\(c_A^2\)}{квадрат коэффициента вариации межприходных интервалов (входящего потока).}
\abbr{\(c_B^2\)}{квадрат коэффициента вариации времени обслуживания (объема работы).}
\abbr{\(\lambda\)}{средняя интенсивность входящего потока заявок.}
\abbr{\(b\)}{средний объём работы заявки; среднее время обслуживания.}
\abbr{\(\mu = 1/b\)}{параметр экспоненциального распределения времени обслуживания в марковском частном случае.}
\abbr{\(\lambda(n)\)}{интенсивность поступления заявок в состоянии \(n\).}
\abbr{\(\lambda_k\)}{интенсивность поступления заявок в состоянии \(k\).}
\abbr{\(\mu(n)\)}{интенсивность обслуживания в состоянии \(n\) в редуцированной модели.}
\abbr{\(\sigma_k\)}{суммарная скорость обслуживания в состоянии \(k\) в ресурсной модели.}
\abbr{\(\alpha_n\)}{коэффициент изменения скорости обслуживания после поступления заявки в state-dependent \(M/G/1\)-модели.}

\abbr{\(T_{\max}\)}{модельное время (длительность одной репликации в имитационной модели).}
\abbr{\(T_{\mathrm{warm}}\)}{период разогрева имитационной модели до выхода на стационарный режим.}
\abbr{\(R_{\mathrm{rep}}\)}{количество независимых репликаций в методе Монте-Карло.}

\abbr{\(\xi(t)\)}{число заявок в системе в момент времени \(t\).}
\abbr{\(\gamma(t)\)}{вектор или матрица занятых ресурсов в момент времени \(t\).}
\abbr{\(\delta(t)\)}{вектор суммарно занятых ресурсов по всем типам.}
\abbr{\(X(t)\)}{случайный процесс, описывающий состояние системы.}
\abbr{\(\pi(n)\)}{стационарная вероятность состояния \(n\).}
\abbr{\(p_0\)}{стационарная вероятность пустой системы.}
\abbr{\(p_k(r_1, \dots, r_k)\)}{стационарное распределение по числу заявок и занятым ими ресурсам.}

\abbr{\(\mathbb{E}[N], \mathbb{E}[\xi]\)}{среднее число заявок в системе.}
\abbr{\(\mathbb{P}\{N = K\}\)}{вероятность полной загрузки системы в редуцированной постановке.}
\abbr{\(P_{\mathrm{loss}}\)}{вероятность потери (блокировки) заявки.}
\abbr{\(\Delta P_{\mathrm{loss}}\)}{абсолютное смещение вероятности отказа (отклонение, выраженное в процентных пунктах).}
\abbr{\(P_{\mathrm{serv\mbox{-}loss}}\)}{вероятность отказа из-за отсутствия свободного прибора.}
\abbr{\(\theta, \Lambda_{\mathrm{eff}}\)}{абсолютная пропускная способность системы (число обслуженных заявок в единицу времени).}